\section{n1. Independent canonical envelope reconstruction on every excellent reduced germ; uniform convergence remains the decisive condition.}
\textit{Supports: G331.}

Independent reconstruction of the canonical envelope argument, relative to the ENTIRE qualified G331. All its source annotations, model restrictions and reading limits remain attached. Let R be excellent local, reduce nilpotents, fix a proper primary ideal a, and write X=V\_a(R), C=C\_geom and D=D\_geom. Every prime support remains. G331 grants the independently defined actual negative effective real Cartier core with nonnegative degree on EVERY contracted integral curve on directed ACTUAL primary blowups, its comparison with ideal C/D, and positive full-fibre descent. It grants no envelope theorem.

For finite continuous signed f set A\_0=R and A\_m=\{g in R: v(g)+m f(v)>=0 for every v in X\}. If M=max(0,-min f), then a\textasciicircum{}N subset A\_m for positive N>=mM. The ultrametric inequality and v(R)>=0 prove idealness. A negative obstacle value excludes every unit, so A\_m is proper primary; if f>=0 everywhere A\_m=R. Multiplicativity gives A\_m A\_n subset A\_(m+n). In a monic ideal-integrality equation for g, v(g)<v(A\_m) makes the leading term uniquely least and finite, a contradiction. Thus A\_m is integrally closed. Define h\_m=-v(A\_m)/m. Its primary Cartier presentation is negative effective and relatively ample; units give zero. Therefore h\_m belongs to the independently geometric C and h\_m<=f.

Let p(f)=sup\{u in D:u<=f\}, initially only a pointwise function. It equals sup of rational primary/unit profiles below f and equals sup\_m h\_m. A rational competitor -v(I)/r<=f has I subset A\_r. Conversely every C competitor is uniformly approximable by rational profiles; subtract a sufficiently small positive rational a-constant, realized by multiplying powered ideals, to obtain competitors converging from below. For u=inf\_i g\_i<=f in an actual decreasing C-net, the open sets \{g\_i<f+epsilon\} cover compact X and are directed under inclusion. One index works everywhere. Then g\_i-epsilon is a C competitor below f and u<=p\_C(f)+epsilon. This proves p\_D=p\_C, including bottom and partial poles. No supremum positivity is inferred here.

Set M\_1=||v(A\_1)|| and H\_N=max\_(1<=r<=N)h\_r in C. For m=kr+s, 0<=s<r, A\_r\textasciicircum{}k A\_1\textasciicircum{}s subset A\_m gives h\_m>=(kr/m)h\_r-sM\_1/m>=h\_r-rM\_1/m because h\_r<=0. Hence H\_N-NM\_1/m<=h\_m<=p. Every degree h\_m converges pointwise to finite p. Since p is a supremum of continuous functions, it is lower semicontinuous. It is continuous iff H\_N converges uniformly iff h\_m converges uniformly: continuity gives the increasing compact strict-error cover for H\_N; the displayed bound transfers uniform convergence to all degrees; the converse follows from uniform limits. When continuous, p belongs to closed geometric C and is the greatest D minorant. This is the precise missing general regularity step.

Subtracting an epsilon a-constant from a D competitor below f makes it a competitor below g whenever f<=g+epsilon. Reversing f and g proves ||p(f)-p(g)||<=||f-g||. In particular uniform obstacle density transfers CE once envelopes for the dense family are actually positive and continuous.

The fixed-obstacle finite-generation equivalence in current413 is also valid at its stated scope. If the canonical algebra is generated by homogeneous g\_i t\textasciicircum{}(d\_i), let L be divisible by all d\_i and J=(g\_i\textasciicircum{}(L/d\_i)). Any degree-m expression has value at least m min\_i v(g\_i)/d\_i. Pure powers attain this bound in degree L, so p=h\_L=-v(J)/L=max\_(r<=max d\_i)h\_r. J is primary or unit: its radical contains A\_1, and the proper case lies in m. The residue bound gives uniform error <=(max d\_i)M\_1/m. Conversely, if p=-v(J)/q, its defining inequalities imply A\_m=\{g:g\textasciicircum{}q in integral-closure(J\textasciicircum{}m)\}. The centred rank-one DVR criterion is applied to each positive-dimensional minimal-component domain and then combined by SH1.1.5 and6.8.3--4. Every such DVR normalizes into X. Monic equations prove the reverse inequalities on every other support. A zero-dimensional minimal component in a reduced local positive-dimensional ring would make m minimal and force dimension zero, so none is lost.

To prove finiteness without assuming a section algebra is finitely generated, put B=R[Jt\textasciicircum{}q]. A primary J contains a regular element. B is reduced excellent with total quotient product K\_j(t\textasciicircum{}q). Excellent/Nagata finiteness of integral closure in the finite product K\_j(t), also when inseparable, supplies a finite B-module. An equation for g\textasciicircum{}q over J\textasciicircum{}m, multiplied by the requisite powers of t\textasciicircum{}(mq), makes g\textasciicircum{}q t\textasciicircum{}(mq) integral over B, then gt\textasciicircum{}m integral by its qth-power equation. The canonical algebra is a B-submodule of this finite module and hence is finite. The unit case is R[t]. A finite maximum of rational profiles is rational by summing suitable powers. This proves canonical FG iff rational attainment iff finite canonical-max attainment, and says nothing about FG density.

For a signed exceptional integral Cartier obstacle f\_D on a normal primary model, pass to its finite normalization Y. The equality A\_m=H\textasciicircum{}0(Y,O(mD)) intersect R follows from exceptional height-one order tests and the normal height-one intersection criterion. At other height-one primes g is regular and D vanishes. Conversely a section is pulled to the valuation ring of Frac(R/p) at each support p: its generic point is outside the exceptional fibre, so the invertible lattice has its nonzero generic trivialization and gives v(g)>=-m f\_D(v). No fraction field of R is mapped through nonzero p. This proves the section formula, not general FG.

For another primary b let d(v)=v(b)>0 and T(v)=v/d(v). Transport f by f\_b(Tv)=f(v)/d(v). All defining inequalities divide by d, so every A\_m and its algebra are literally unchanged; h\_m and p divide by d. This gives the exact gauge cocycle and transports uniform convergence since d and 1/d are bounded. For normalization S, surjective restriction gives A\_m\textasciicircum{}R=A\_m\textasciicircum{}S intersect R. A continuous source envelope descends to the downstairs envelope if it is constant on EVERY actual fibre, including the infinite-conductor quotient fibres, by G331's positive C descent and pullback of D competitors. Obstacle compatibility alone does not prove this constancy. A decreasing family of continuous obstacles with C envelopes has the D envelope equal to the infimum of those envelopes: it is an actual decreasing C-net, lies below the limiting obstacle, and dominates every competitor.

On an excellent reduced curve the normalization centres are DVRs. A regular conductor element has finite positive order at every branch, so normalized boundary fibres are singletons. Every real nonpositive slope is an actual Cartier evaluation on a DVR, with no contracted-curve test. Semilocal gluing and G331 positive descent give the finite nonpositive orthant C; coordinate truncations give the extended orthant D. The geometric envelope is min(f,0), exactly the ideal envelope. Literal classical/CLD multibranch radial extensions remain distinct, retaining the mixed-zero obstruction. A DVR obstacle -alpha, alpha irrational positive, has a real Cartier continuous envelope but A\_m=(x\textasciicircum{}ceil(m alpha)) is not FG, since finite homogeneous generation would force a rational minimum order/degree slope. Non-FG is not a CE failure. Zero-dimensional reduction has empty boundary; nilpotents are invisible.

These arguments reconstruct the used current413 canonical results without using its verdict or compiler record as a theorem. Complete immediate SH units and both formula pages were freshly read/viewed; their named deeper foundations retain the disclosed limits. Artifacts: n1.txt, expanded n1.tex, P12.3-surface-envelopes.tex, proof-companion.md. Rerun: python build\_delivery.py; python check\_delivery.py. Execution is serialization evidence only.

\subsection*{Expanded canonical convergence and gauge formulae}
The full sequence is not presumed monotone. For $m=kr+s$, $0\le s<r$,
\[
 A_r^k A_1^s\subset A_m,\qquad
 H_N-\frac{NM_1}{m}\le h_m\le p(f),\qquad
 p(f)=\sup_m h_m.
\]
The compact increasing-error-cover argument applies to $H_N$ and gives
\[
 p(f)\in C(X)\ \Longleftrightarrow\ H_N\to p(f)\text{ uniformly}
 \ \Longleftrightarrow\ h_m\to p(f)\text{ uniformly}.
\]
Under $T(v)=v/v(\mathfrak b)$ all canonical ideals are unchanged;
only evaluations change by the bounded positive factor $v(\mathfrak b)^{-1}$.
Rational envelope attainment and canonical finite generation are equivalent
for the fixed obstacle; neither is a universal regularity hypothesis.
